\section{Finding Authorities by Subject}
\label{sec:retrieval}
\subsection{Retrievers}
We index the title and headnote of every judgment in two ways. \emph{BM25} \cite{robertson2009} scores unigrams and bigrams (English stop words removed). \emph{Dense} retrieval embeds each headnote with the open multilingual model bge-m3 \cite{chen2024m3}, run locally, and scores by cosine similarity with the embedded query. The \emph{hybrid} retriever takes the union of the top 300 judgments of each and adds their scores after normalising each to zero mean and unit variance over the candidates. A \emph{citation prior} adds the normalised $\log(1+c)$, where $c$ is the number of judgments in the graph that cite the candidate. We use the raw count; recursive measures such as PageRank \cite{page1999} are an alternative that we did not test. Every combination is a weighted sum of these three scores; the weights are chosen on a development split and never on the evaluation data.

\subsection{Evaluation without human labels}
The graph supplies relevance judgments without annotation: the judgments a decision cites are the authorities its bench considered, the standard setting of citation recommendation \cite{farber2020} and of Indian prior-case retrieval \cite{joshi2023ucreat}. For each held-out judgment that cites at least \RecMinTargets{} earlier judgments, the query is its headnote with every reporter citation, every ``X v.\ Y'' case name and every ``\ldots's case'' reference removed, so the query states the legal issues without naming the precedents. The targets are the judgments it cites. Only judgments decided before it are candidates, and the citation prior counts only citations made before it, so no information from the future is used. We choose weights on \RecNDev{} development judgments and report recall@10, recall@50, the share of queries with at least one target in the top 10 (hit@10) and mean reciprocal rank on \RecNTest{} other test judgments, with bootstrap 95\% intervals. We also score the 40 T4 questions against their reference judgments, with the weights from the development split.
